% Teaser: Progressive Hamiltonian Landscape Learning under Local Sensing
\begin{figure}[t]
\centering
\begin{tikzpicture}[
  font=\small,
  >=Latex,
  box/.style={draw, rounded corners=2pt, very thick, inner sep=1pt, fill=black!1},
  rowtag/.style={draw, rounded corners=6pt, fill=black!6, inner xsep=6pt, inner ysep=2pt,
                 font=\scriptsize\bfseries, align=center},
  header/.style={draw, rounded corners=3pt, fill=black!10, very thick,
                 inner xsep=6pt, inner ysep=3pt, font=\scriptsize\bfseries, align=center},
  legendbox/.style={draw, rounded corners=2pt, thick, fill=white, inner sep=5pt},
  arrow/.style={->, very thick},
  accent/.style={draw=RoyalBlue, very thick},
  accent2/.style={draw=Crimson,  very thick},
  accent3/.style={draw=ForestGreen, very thick},
  alabel/.style={font=\scriptsize, fill=white, rounded corners=2pt, inner sep=2pt}
]

% ------------------- 2-column grid: Physical | Energy -------------------
\matrix (G) [matrix,
             row sep=10mm,
             column sep=38mm,
             nodes={anchor=center}] at (0,0) {

  \node[box] (phys0) {\includegraphics[width=7.0cm,height=2.35cm,keepaspectratio]{figs/phys_try1.png}}; &
  \node[box] (eng0)  {\includegraphics[width=7.0cm,height=2.35cm,keepaspectratio]{figs/energy_try1.png}}; \\

  \node[box] (phys1) {\includegraphics[width=7.0cm,height=2.35cm,keepaspectratio]{figs/phys_try2.png}}; &
  \node[box] (eng1)  {\includegraphics[width=7.0cm,height=2.35cm,keepaspectratio]{figs/energy_try2.png}}; \\

  \node[box] (phys2) {\includegraphics[width=7.0cm,height=2.35cm,keepaspectratio]{figs/phys_try3.png}}; &
  \node[box] (eng2)  {\includegraphics[width=7.0cm,height=2.35cm,keepaspectratio]{figs/energy_try3.png}}; \\
};

\node[header, text width=4.5cm, anchor=south] (colL) at ($(phys0.north)+(0,7mm)$)
{Physical space\\(unknown $\rightarrow$ revealed by local sensing)};

\node[header, text width=4.0cm, anchor=south] (colR) at ($(eng0.north)+(0,7mm)$)
{Hamiltonian energy terrain\\(learned progressively)};

% ------------------- Title -------------------
\node[
  font=\small\bfseries,
  anchor=south,
  text width=13cm,
  align=center
]
at ($(colL.north)!0.5!(colR.north)+(0,2.5mm)$)
{Progressive Hamiltonian Landscape Learning from Local Sensing\\
along Hamiltonian geodesic pathways};

% ------------------- Vertical progression arrows -------------------
\draw[arrow] (phys0.south) -- (phys1.north);
\draw[arrow] (phys1.south) -- (phys2.north);
\draw[arrow] (eng0.south) -- (eng1.north);
\draw[arrow] (eng1.south) -- (eng2.north);

% ------------------- Cross-column arrows + boxed labels -------------------
\draw[arrow, accent] (phys0.east) -- (eng0.west)
  node[midway, above=2mm, alabel] {scan local env $\Rightarrow$ build initial $H_{\tau}$};

\draw[arrow, accent2] (phys1.east) -- (eng1.west)
  node[midway, above=2mm, alabel] {Try 1 feedback $\Rightarrow$ update $H_{\tau}$};

\draw[arrow, accent3] (phys2.east) -- (eng2.west)
  node[midway, above=2mm, alabel] {Try 2 feedback $\Rightarrow$ refine $H_{\tau}$};

% ------------------- Bottom idea box -------------------
\node[draw, rounded corners=2pt, thick, fill=black!1,
      inner sep=6pt, text width=13cm, align=center, anchor=north]
  at ($(G.south)+(0,-11mm)$) {
  \textbf{Idea:} In an unknown environment, the agent repeatedly senses a \emph{local} neighborhood,
  attempts to reach the stage goal, and incrementally learns a Hamiltonian energy landscape whose valleys guide future attempts under partial observability.
};

\end{tikzpicture}
\caption{\textbf{Conceptual overview: progressive energy-landscape shaping from local sensing.}
\textit{Left}: the physical workspace is initially unknown and is revealed only through local sensing around the agent, producing a partial view of nearby obstacles and a stage goal.
\textit{Right}: from this partial context, the agent builds an initial internal “energy terrain” and then repeatedly attempts a short plan toward the stage goal; each attempt produces simple feedback (e.g., collisions/clearance and progress), which is used to refine the terrain. Over successive refinements, the terrain becomes smoother and its low-energy valleys align with traversable corridors, guiding future attempts even under partial observability.}

\label{fig:teaser_progressive_landscape}
\end{figure}
